\SourceLeaf{212}
\textit{utilissima regula} nennt es \textsc{Cotes}) bewogen, diese Vorschriften weiter und bis auf
den Fall von 11 Ordinaten auszudehnen. Immer erscheint so der verlangte
Flächenraum in der Gestalt des Products der Basis, oder der Entfernung der
äussersten Ordinaten, in die Summe der durch bestimmte Zahlcoëfficienten multiplicirten
Ordinaten, und zwar haben zwei gleich weit vom Anfang und Ende
abliegende Ordinaten allemal gleiche Coëfficienten. Diese Quadraturcoëfficienten
bis zu dem Fall von 11 Ordinaten gibt \textsc{Cotes} am Schluss der Abhandlung \textit{de methodo differentiali},
welche einen Theil der \textit{Harmonia mensurarum} ausmacht, ohne
sich über das Verfahren, wodurch er sie berechnet hat, weiter zu erklären. Vielleicht
hat man es dieser anspruchlosen Kürze, womit bloss das Endresultat dargestellt
ist, zuzuschreiben, dass diese schöne und zweckmässige Methode von den
Analysten weniger gekannt und benutzt zu sein scheint, als sie es verdient.

Bei dieser Methode liegt durchaus die Voraussetzung gleicher Abstände zwischen
den Ordinaten zum Grunde. Allerdings scheint beim ersten Anblick diese
Voraussetzung am einfachsten und natürlichsten zu sein, und es war noch nicht
in Frage gekommen, ob es nicht demungeachtet noch vortheilhafter sein könne,
Ordinaten in ungleichen Abständen zum Grunde zu legen. Um diese Frage zu
entscheiden, musste zuerst die Theorie der Quadraturcoëfficienten in unbeschränkter
Allgemeinheit entwickelt, und der Grad der Genauigkeit des Resultats bestimmt
werden. Es zeigte sich, dass die Bedingungen, wovon dieser Grad der
Genauigkeit abhängt, von der Art sind, dass man dieselbe durch zweckmässig
gewählte Ordinaten in ungleichen Abständen allerdings verdoppeln kann, so dass
man mit einer beliebigen Anzahl gehörig gewählter Ordinaten eben so weit reicht,
als mit der doppelten Anzahl von Ordinaten in gleichen Abständen. Diese Untersuchungen,
nebst der vollständigen Theorie der zweckmässigsten Auswahl der
Ordinaten, der dabei anzuwendenden Quadraturcoëfficienten und der Bestimmung
des Grades der Genauigkeit, welchen dieses Verfahren gewährt, machen den
Hauptinhalt der vorliegenden Abhandlung aus.

Aus der kurzen Entwickelung der Theorie der \textsc{Cotes}ischen Quadraturcoëfficienten,
welche der Verf. vorausschicken zu müssen glaubte, berühren wir hier
nur dasjenige, was den Grad der Genauigkeit betrifft, welchen die dadurch gefundenen
genäherten Integrale haben. Vor allen muss hier bemerkt werden, dass
die Anwendbarkeit dieser Methode, eben so wie das Interpoliren, auf der Voraussetzung
beruhe, dass die Ordinaten innerhalb des zu quadrirenden Raumes
\SourceLeaf{213}
sich durch eine convergirende Reihe darstellen lassen. Es sei \(x\) die Abscisse, \(y\)
die Ordinate, und das Integral \(\int y\,\mathrm{d}x\) werde von \(x=g\) bis \(x=h\) verlangt.
Man führe statt \(x\) eine andere veränderliche Grösse ein, indem man etwa
\(x=g+(h-g)t\), oder auch \(x=\frac12(g+h)+\frac12(h-g)u\) setzt. Hier muss also
\(y\) sich durch Reihen wie
\[
\alpha+\alpha't+\alpha''tt+\alpha'''t^3+\text{etc.}
\]
oder
\[
\beta+\beta'u+\beta''uu+\beta'''u^3+\text{etc.}
\]
darstellen lassen, die convergiren, jene, wenigstens so lange \(t\), diese so lange \(u\)
nicht grösser wird als 1. Man mag daher Kürze wegen den Coëfficienten \(\alpha'\) und \(\beta'\)
die Ordnung 1, den Coëfficienten \(\alpha''\) und \(\beta''\) die Ordnung 2 u.\,s.\,w. beilegen. Diess
vorausgesetzt, wird gezeigt, dass die Fehler, denen man sich bei der \textsc{Cotes}ischen
Methode aussetzt, zwar immer von einer höhern Ordnung werden, je grösser die
Anzahl der zum Grunde gelegten Werthe von \(y\) ist, jedoch so, dass eine ungerade
Anzahl und die zunächst grössere gerade Anzahl immer Fehler von einerlei
Ordnung hervorbringen. So ist für drei Ordinaten der Fehler sehr nahe
\(=\frac{1}{120}(h-g)\alpha''''\), für vier Ordinaten nahe \(=\frac{1}{270}(h-g)\alpha''''\); sodann für
fünf Ordinaten nahe \(=\frac{1}{2688}(h-g)\alpha^{\mathrm{VI}}\) und für sechs Ordinaten nahe
\(=\frac{11}{52500}(h-g)\alpha^{\mathrm{VI}}\) u.\,s.\,f. Man sieht hieraus, dass es im Allgemeinen vortheilhaft
sein wird, bei Anwendung der \textsc{Cotes}ischen Methode eine ungerade Anzahl von
Ordinaten zu benutzen.

Der Verf. geht hierauf zu der allgemeinen Untersuchung über, wo die Einschränkung,
dass die Ordinaten gleiche Abstände von einander haben, wegfällt.
Sind hier \(A,A',A''\) u.\,s.\,w. die Werthe von \(y\). die entsprechenden Werthe von \(t\)
hingegen \(a,a',a''\) u.\,s.\,w., oder \(b,b',b''\) u.\,s.\,w. die entsprechenden Werthe von \(u\), und
ihre Anzahl \(n+1\), so wird das genäherte Integral wiederum die Gestalt haben
\[
(h-g)(RA+R'A'+R''A''+\text{etc.})
\]
wo \(R,R',R''\) u.\,s.\,w. Zahlcoëfficienten sind, die unabhängig von der Function \(y\)
bloss durch \(a,a',a''\) u.\,s.\,w., oder durch \(b,b',b''\) u.\,s.\,w. bestimmt werden. Die Untersuchungen
des Verf. geben für diese Bestimmung folgendes Resultat. Es sei
\[
T=(t-a)(t-a')(t-a'')\,\ldotp\,\ldotp\,\ldotp\,\ldotp
\]
\SourceLeaf{214}
Aus der Multiplication dieser ganzen Function von \(t\), welche auf die Ordnung
\(n+1\) steigt, in die unendliche Reihe, welche den Logarithmen von \(\frac{t}{t-1}\) vorstellt,
nemlich
\[
t^{-1}+\frac12t^{-2}+\frac13t^{-3}+\text{etc.}
\]
ergebe sich das Product \(T'+T''\), so dass \(T'\) die darin enthaltene ganze Function
von \(t\) bezeichnet, so wie \(T''\) die übrige mit negativen Potenzen von \(t\) ins Unendliche
fortlaufende Reihe. Diess vorausgesetzt, ergeben sich die Quadraturcoëfficienten
\(R,R',R''\) u.\,s.\,w., wenn man in \(\frac{T'\mathrm{d}t}{\mathrm{d}T}\) für \(t\) der Reihe nach die Werthe
\(a,a',a''\) u.\,s.\,w. substituirt. Auf eine ähnliche und noch etwas bequemere Art
leitet man jene Coëfficienten aus \(b,b',b''\) u.\,s.\,w. ab, indem man die Function
\[
(u-b)(u-b')(u-b'')\,\ldotp\,\ldotp\,\ldotp\,\ldotp
\]
durch \(U\), ihr Product in die unendliche Reihe
\[
u^{-1}+\frac13u^{-3}+\frac15u^{-5}+\text{etc.}
\]
durch \(U'+U''\) bezeichnet (so dass \(U'\) die darin enthaltene ganze Function von
\(u\) vorstellt), und dann in \(\frac{U'\mathrm{d}u}{\mathrm{d}U}\) für \(u\) der Reihe nach die Werthe \(b,b',b''\) u.\,s.\,w.
substituirt. Statt der gebrochenen Functionen \(\frac{T'\mathrm{d}t}{\mathrm{d}T},\frac{U'\mathrm{d}u}{\mathrm{d}U}\) lassen sich auch ganze
Functionen von \(t\) und \(u\) finden, welche die Stelle von jenen vertreten können,
und für deren Bestimmung der Verf. eine allgemeine Methode entwickelt.

Der Grad der Genauigkeit der Integrationsformel hängt nun von der Beschaffenheit
der Reihe \(T''\) oder \(U''\) ab. Im Allgemeinen ist der Fehler zwar von
der Ordnung \(n+1\); allein wenn von den ersten Gliedern jener Reihen einige
ausfallen, so wird der Fehler von einer höhern Ordnung, so dass wenn \(T''\) erst
mit der Potenz \(t^{-m}\) oder \(U''\) mit der Potenz \(u^{-m}\) anfängt, der Fehler von der
Ordnung \(n+m\) wird.

Hieraus ergab sich nun, dass in so fern die Werthe \(a,a',a''\) u.\,s.\,w., oder
\(b,b',b''\) u.\,s.\,w. willkürlich gewählt werden können, diese sich so bestimmen lassen
müssen, dass die ersten \(n+1\) Glieder von \(T''\) oder \(U''\) wirklich ausfallen;
wovon die Folge sein wird, dass der Fehler der Integrationsformel auf die Ordnung
\(2n+2\) kommt. Die Untersuchung schreitet demnach zu der Bestimmung
derjenigen Functionen \(T\) und \(U\), für jeden Werth von \(n\), fort, wodurch der
angegebenen Bedingung Genüge geleistet wird. Der beschränkte Raum erlaubt
\SourceLeaf{215}
uns nicht, in das Einzelne dieser Untersuchung hier einzugehen: wir bemerken
also hier nur, dass diese Functionen ein sehr einfaches Fortschreitungsgesetz
befolgen, und in genauem Zusammenhange stehen mit der Entwicklung der Reihen
\[
\begin{aligned}
&t^{-1}+\tfrac{1}{2}t^{-2}+\tfrac{1}{3}t^{-3}+\tfrac{1}{4}t^{-4}+\mathit{etc.}\\
&u^{-1}+\tfrac{1}{3}u^{-3}+\tfrac{1}{5}u^{-5}+\tfrac{1}{7}u^{-7}+\mathit{etc.}
\end{aligned}
\]
in continuirliche Brüche, die der Verf. in einer frühern Abhandlung
[\textit{Disquisitiones generales circa seriem infinitam}
\(1+\frac{\alpha.\beta}{1.\gamma}x+\mathit{etc.}\)] gegeben hat.\,---
Offenbar gibt demnächst die Auflösung der Gleichung \(T=0\) oder \(U=0\)
die Werthe von \(a,a',a''\) u.\,s.\,w. oder \(b,b',b''\) u.\,s.\,w.,
und die Werthe der Quadraturcoëfficienten werden nach den allgemeinen Regeln
bestimmt, die in diesem Falle noch besondere Vereinfachungen vertragen.
Uebrigens werden allerdings in den meisten Fällen sowohl die Werthe von
\(a,a',a''\) u.\,s.\,w. als die Quadraturcoëfficienten Irrationalgrössen.
Diess ist indess an sich sehr gleichgültig, sobald nur ihre numerischen
Werthe ein für allemal mit einem angemessenen Grad von Genauigkeit berechnet
sind. Ist diess der Fall, so wird die Anwendung dieser Methode auf irgend
eine Anzahl von Ordinaten wenig oder gar nicht mehr Mühe machen, als die
Anwendung der \textsc{Cotes}ischen Methode auf eine eben so grosse Anzahl,
da hingegen letztere auf eine doppelt so grosse Anzahl angewandt werden
müsste, um ungefähr dieselbe Genauigkeit des Resultats zu geben, wie erstere.

Um für die Anwendung dieser neuen Methode nichts zu wünschen übrig zu
lassen, hat der Verf. noch die numerischen Werthe von \(a,a',a''\)
u.\,s.\,w., so wie von \(R,R',R''\) u.\,s.\,w., auf 16 Decimalen
berechnet, mitgetheilt, zugleich mit den \textsc{Brigg}ischen Logarithmen
der letztern auf 10 Decimalen, alles bis zu dem Fall von sieben Ordinaten.
In diesem letzten Fall wird der Fehler der Integrationsformel nahe
\(=\frac{1}{176679360}\alpha^{\mathrm{XIV}}\), woraus man abnehmen
kann, dass in den meisten in der Ausübung vorkommenden Fällen schon eine
geringere Anzahl zureichen wird. Um die Anwendung der Vorschriften und ihre
verhältnissmässige Schärfe noch mehr zu versinnlichen, ist als Beispiel die
Berechnung von \(\int\frac{\mathrm{d}x}{\log x}\) von \(x=100000\)
bis \(x=200000\) beigefügt, wo schon bei der Anwendung von vier Werthen
der Fehler nur \(\frac{1}{56000000}\) des Ganzen ist, und bei einer
grössern Anzahl sich in den unvermeidlichen Fehlern verliert, die selbst
die grössern Logarithmentafeln noch übrig lassen.
\par\begin{center}\rule{0.22\linewidth}{0.3pt}\end{center}
